% DRAFT: paper §6.x — E7 b=2 mechanism test, combined Phase 1 / 2a / 2b.
% Closing entry 2026-06-14. Phase 1 confirms the b=2/TIES covariance;
% Phase 2b cleanly REJECTS the "implicit-sparsity" explanation.

\subsection{What drives the $b=2$ dip}
\label{subsec:exp-b2-mechanism}

The multi-seed merge matrix of \S\ref{subsec:exp-multiseed} shows
that $2$-bit task-vector quantization (TVQ) often outperforms higher
bit budgets at worst-task NLL excess --- the ``less-is-more'' dip
that survives across all four base models. The original paper
\S\ref{sec:discussion} conjectured an \emph{implicit-TIES}
explanation: $b=2$ quantization implicitly induces sparsity, which
plays the same role as the explicit magnitude trimming TIES
performs. We pre-registered two falsifiable predictions of that
hypothesis and tested both.

\paragraph{Prediction 2 (covariance with TIES). \textsc{Confirmed.}}
If $b=2$ TVQ and TIES share a mechanism, their per-task excess
patterns should covary across the merge matrix. Pairing each
\textsc{(model, task, seed)} cell's
$\text{dip}_t \coloneqq \Delta L_t^{(b=4)} - \Delta L_t^{(b=2)}$
with that cell's
$\text{ties\_win}_t \coloneqq \Delta L_t^{(\mathrm{TA})}
- \Delta L_t^{(\mathrm{TIES})}$,
the Spearman rank correlation across the 32 cells is
$\rho = +0.96$ (Pearson $r = +0.95$), well above the pre-registered
$0.7$ threshold. Per-model the pattern holds 4/4 (Llama-3.1-8B
$\rho=0.98$, Mistral $0.95$, Qwen $0.95$, Yi $0.81$); per-task it
holds 3/4 (Alpaca $0.93$, Magicoder $0.95$, Translation $0.83$;
GSM8K is a within-task null, where both interventions reliably
help but the variation does not covary). \emph{So $b=2$ TVQ and
TIES are doing the same thing to per-task error, by some shared
mechanism.}

\paragraph{Prediction 1 (sparsity is the shared mechanism). \textsc{Rejected.}}
$b=2$ quantization induces substantial implicit sparsity: across
the 16 task-adapter deltas, the post-quantization tensor places
$50\%$ of coordinates in the bucket containing zero, and the
quantization error exceeds the original magnitude for $\sim 87\%$
of coordinates. To test whether this sparsity is what drives the
dip, we constructed an explicit magnitude-pruning merge:
per-layer, keep top-$\rho$ fraction by $\lvert\Delta\rvert$, zero
rest, sum without rescaling, SVD-truncate to rank-$r$. We evaluated
two pruning densities matching the two implicit-sparsity measurements
($\rho = 0.50$ and $\rho = 0.13$) on the four base models.

\begin{table}[t]
\centering
\caption{Worst-task NLL excess of explicit magnitude-pruning at the
two densities that match $b=2$ TVQ's implicit sparsity, with
$b=2$ and $b=4$ TVQ for reference. Pruning at the matched
sparsity \emph{does not} reproduce the $b=2$ dip; in 7 of 8 cells
it underperforms even $b=4$.}
\label{tab:e7-prune-vs-tvq}
\begin{tabular}{lrrrr}
\hline
Model & magprune $\rho=0.50$ & magprune $\rho=0.13$ & TVQ $b=2$ & TVQ $b=4$ \\
\hline
Llama-3.1-8B    & $0.225$ & $0.288$ & $0.108$ & $0.217$ \\
Mistral-7B-v0.3 & $0.155$ & $0.251$ & $0.058$ & $0.138$ \\
Qwen-2.5-7B     & $0.118$ & $0.229$ & $0.019$ & $0.105$ \\
Yi-1.5-9B       & $0.088$ & $0.151$ & $0.059$ & $0.098$ \\
\hline
\end{tabular}
\end{table}

In zero of four models does $\rho = 0.50$ pruning hit the $b=2$
dip; only Yi shows a partial result ($0.088$ between $b=2$ at
$0.059$ and $b=4$ at $0.098$). At $\rho = 0.13$ (matching the
fraction of coordinates whose quantization error exceeds their
magnitude) pruning underperforms every baseline. \emph{Implicit
sparsity is not the mechanism behind the $b=2$ dip.}

\paragraph{Sharpened hypothesis.}
The covariance with TIES is real (Prediction 2), but it is
not mediated by implicit sparsity (Prediction 1). What \emph{is}
shared between min-max scalar quantization at $b=2$ and
TIES's trim$+$sign-elect$+$disjoint-mean? Two candidates the
data leave open:
\begin{enumerate}
\item \emph{Uniform shrinkage of large coefficients.} A min-max
quantizer with $3$ levels maps every coordinate to one of three
bucket centers, so the largest-magnitude coefficients are
pulled \emph{toward} the smaller surviving bucket. TIES does
not have this effect explicitly, but its sign-election step
under-counts coefficients whose elected and per-task signs
disagree, producing an effective shrinkage of the sum's
extremes. Both interventions therefore compress the
distribution of merged coefficients, even though the
mechanisms by which they do so look superficially unrelated.
\item \emph{Per-tensor (rather than per-coordinate) calibration.}
Both methods normalize at the tensor level --- TVQ via the
per-tensor $\min/\max$, TIES via the per-tensor magnitude
threshold and elected signs. Magnitude pruning at a global
density is also per-tensor but lacks the per-tensor
\emph{distributional} re-anchoring that quantization and
sign-election provide.
\end{enumerate}
Disentangling these is a follow-up question for a focused
mechanism study (we sketch the design in
\S\ref{sec:future}); the present finding is that the $b=2$ dip
is real, covaries with TIES at the per-task level, and is
\emph{not} explained by implicit sparsity. The implicit-TIES
hypothesis as originally stated is refined: $b=2$ and TIES
behave alike because they both impose tensor-level rather than
coordinate-level structure on the merge, not because $b=2$
trims small coefficients.

%%% END DRAFT %%%
